\subsection{复线性表示}

在本小节中，我们假设 K 是复数域 C。

设 \((M,\pi)\) 是 G 的一个线性表示。若 M 上的埃尔米特型 \(\Phi\) 满足

\[
\Phi (\pi (g) x, \pi (g) x ^ {\prime}) = \Phi (x, x ^ {\prime})\tag{48}
\]

（对所有 \(x, x' \in M\) 与每个 \(g \in G\)），则称它在 G 下不变。这也意味着，对每个 \(g \in G\)，M 的自同构 \(\pi(g)\) 关于 \(\Phi\) 是酉的。

命题 11. —— 设 \((M,\pi)\) 是 G 的一个有限维线性表示。

\begin{enumerate}
\def\labelenumi{\alph{enumi})}
\item
  在 M 上存在一个正的、分离的且在 G 下不变的埃尔米特型。
\item
  假设表示 \(\pi\) 是单表示。若 \(\Phi\) 与 \(\Psi\) 是 M 上非零的且在 G 下不变的埃尔米特型，则存在实数 a 使得 \(\Psi = a\Phi\)。
\end{enumerate}

在向量空间 M 上选取一个分离的正埃尔米特型，并把它记作 \(\Phi_{0}\)。我们通过下面的式子定义一个在 G 下不变的分离的正埃尔米特型 \(\Phi\)：

\[
\Phi (x, x ^ {\prime}) = \sum_ {g \in \mathrm{G}} \Phi_ {0} (\pi (g) x, \pi (g) x ^ {\prime})\tag{49}
\]

（对 \(x, x' \in M\)）。

设 \(\Psi\) 是 M 上的一个埃尔米特型；存在 M 的唯一的自同态 A，使得 \(\Psi(x, x') = \Phi(x, Ax')\)（对 M 中的 \(x, x'\)）。此外，若 \(\Psi\) 在 G 下不变，则自同态 A 与自同构 \(\pi(g)\)（对 \(g \in G\)）交换。若表示 \(\pi\) 是单表示，则由舒尔引理（VIII, 第 47 页，定理 1），A 是一个位似，从而存在复数 \(a\) 使得 \(\Psi = a\Phi\)。由于 \(\Phi\) 与 \(\Psi\) 都是埃尔米特型且 \(\Phi\) 非零，\(a\) 是实数。命题得证。

我们赋予向量空间 \(\mathbf{C}[\mathrm{G}]\)（\(\mathrm{G}\) 上复函数所成的向量空间）一个希尔伯特空间结构，其内积由

\[
\langle f | f ^ {\prime} \rangle_ {\mathrm{G}} = | \mathrm{G} | ^ {- 1} \sum_ {g \in \mathrm{G}} \overline {{f (g)}} f ^ {\prime} (g).\tag{50}
\]

给出。对任意函数 \(f \in C[G]\)，用 \(f^{*}\) 表示由

\[
f ^ {*} (g) = \overline {{f (g ^ {- 1})}}\tag{51}
\]

（对 \(g \in \mathbf{G}\)）定义的函数；映射 \(f \mapsto f^*\) 是 \(\mathbf{C}[\mathbf{G}]\) 的一个半线性对合。我们还有

\[
\langle f | f ^ {\prime} \rangle_ {\mathrm{G}} = \langle f ^ {*}, f ^ {\prime} \rangle_ {\mathrm{G}}\tag{52}
\]

（对 \(f, f' \in \mathbf{C}[\mathrm{G}]\)），这里采用 VIII, 第 410 页公式 (28) 的记号。因此我们有

\[
\langle f | f ^ {\prime} \rangle_ {\mathrm{G}} = | \mathrm{G} | ^ {- 2} \tau (f ^ {*} f ^ {\prime}).\tag{53}
\]

设 \((\mathrm{M},\pi)\) 是 G 的一个有限维线性表示。我们赋予向量空间 M 一个使诸自同态 \(\pi(g)\) 为酉自同构的希尔伯特空间结构（命题 11）。若用 \(A^{*}\) 表示 M 的自同态 A 关于这个结构的伴随，则我们有 \(\mathrm{Tr}(A^{*})=\overline{\mathrm{Tr}(A)}\)。对每个 \(g\in G\)，我们有 \(\pi(g^{-1})=\pi(g)^{*}\)，从而 \(\chi_{\pi}(g^{-1})=\overline{\chi_{\pi}(g)}\)；换言之，我们有 \(\chi_{\pi}=\chi_{\pi}^{*}\)。于是特征标的正交关系（VIII, 第 410 页，命题 4）取如下形式

\[
\langle \chi_ {\lambda} | \chi_ {\mu} \rangle_ {\mathrm{G}} = \delta_ {\lambda \mu}\tag{54}
\]

（对 \(\lambda, \mu \in \widehat{\mathbf{G}}\)）。它表达了这样一个事实：特征标族 \((\chi_{\lambda})_{\lambda \in \widehat{\mathbf{G}}}\)——\(\mathbf{G}\) 的单表示的特征标族——是由中心函数组成的希尔伯特空间 \(\mathrm{Z}(\mathbf{C}[\mathrm{G}])\) 的一组规范正交基。

设 \(\pi\) 与 \(\pi'\) 是 G 的有限维线性表示。我们有关系式 \(\langle \chi_{\pi}|\chi_{\pi'}\rangle_{\mathrm{G}} = \dim_{\mathbf{C}}\mathrm{Hom}_{\mathrm{G}}(\pi, \pi')\)（VIII, 第 410 页，推论）。表示 \(\pi\) 不可约，当且仅当 \(\langle \chi_{\pi}|\chi_{\pi}\rangle_{\mathrm{G}} = 1\)。

对每个元素 \(\lambda\)（\(\widehat{G}\) 的元素），我们赋予向量空间 \(V_{\lambda}\) 一个使诸自同构 \(\pi_{\lambda}(g)\) 为酉自同构的希尔伯特空间结构。用 \(\langle v|v'\rangle_{\lambda}\) 表示两个元素 \(v, v'\)（\(V_{\lambda}\) 的元素）的内积，用 \(u^{*}\) 表示自同态 \(u\)（\(V_{\lambda}\) 的自同态）的伴随。设 \(A = (A_{\lambda})_{\lambda \in \widehat{G}}\) 与 \(A' = (A_{\lambda}')_{\lambda \in \widehat{G}}\) 是 \(\mathrm{F}(\widehat{\mathrm{G}})\) 的元素。我们记 \(\mathrm{A}^* = (\mathrm{A}_\lambda^*)_{\lambda \in \widehat{\mathrm{G}}}\)。我们有 \(\overline{\mathcal{F}} (a^{*}) = (\overline{\mathcal{F}} (a))^*\)（对每个元素 \(a\)（\(\mathbf{C}[\mathrm{G}]\) 中的元素））。令

\[
\langle \mathrm{A} | \mathrm{A} ^ {\prime} \rangle_ {\widehat {\mathrm{G}}} = | \mathrm{G} | ^ {- 2} \widehat {\tau} (\mathrm{A} ^ {*} \mathrm{A} ^ {\prime}).\tag{55}
\]

由 VIII, 第 407 页的公式 (10)，我们有

\[
\langle \mathrm{A} | \mathrm{A} ^ {\prime} \rangle_ {\widehat {\mathrm{G}}} = \frac {1}{| \mathrm{G} | ^ {2}} \sum_ {\lambda \in \widehat {\mathrm{G}}} d _ {\lambda} \operatorname{Tr} (\mathrm{A} _ {\lambda} ^ {*} \mathrm{A} _ {\lambda} ^ {\prime}).\tag{56}
\]

由于 \(\hat{\tau}\circ\overline{{\mathcal{F}}}=\tau\)，公式 (53) 与 (55) 蕴含映射 \(\overline{{\mathcal{F}}}\) 是从 C{[}G{]} 到 F( \(\widehat{G}\) ) 的希尔伯特空间同构。

舒尔正交关系（VIII, 第 410 页）可以用希尔伯特内积重新表述。于是关系式 (21) 与 (24) 给出下述断言。对 \(\lambda \in \widehat{\mathbf{G}}\) 与 \(x, x', y, y'\)（\(\mathrm{V}_{\lambda}\) 中的元素），我们有

\[
| \mathrm{G} | ^ {- 1} \sum_ {g \in \mathrm{G}} \overline {{\langle x | \pi_ {\lambda} (g) x ^ {\prime} \rangle}} _ {\lambda} \langle y | \pi_ {\lambda} (g) y ^ {\prime} \rangle_ {\lambda} = d _ {\lambda} ^ {- 1} \overline {{\langle x | y \rangle}} _ {\lambda} \langle x ^ {\prime} | y ^ {\prime} \rangle_ {\lambda}.\tag{57}
\]

若 \(\lambda\) 与 \(\mu\) 是 \(\widehat{\mathrm{G}}\) 的两个不同元素，则对 \(x, x'\)（\(\mathrm{V}_{\lambda}\) 中的元素）与 \(y, y'\)（\(\mathrm{V}_{\mu}\) 中的元素），我们有

\[
\sum_ {g \in \mathrm{G}} \overline {{\langle x | \pi_ {\lambda} (g) x ^ {\prime} \rangle_ {\lambda}}} \langle y | \pi_ {\mu} (g) y ^ {\prime} \rangle_ {\mu} = 0.\tag{58}
\]

对每个 \(\lambda \in \widehat{\mathbf{G}}\)，取定规范正交基 \((e_{\lambda,i})_{1\leqslant i\leqslant d_{\lambda}}\)（\(\mathrm{V}_{\lambda}\) 的一组规范正交基）。对任意 \(g\in \mathbf{G}\)，用 \((\pi_{ij}^{\lambda}(g))\) 表示自同态 \(\pi_{\lambda}(g)\)（\(\mathrm{V}_{\lambda}\) 的自同态）关于这组基的矩阵；我们有

\[
\pi_ {i j} ^ {\lambda} (g) = \langle e _ {\lambda , i} | \pi_ {\lambda} (g) e _ {\lambda , j} \rangle_ {\lambda}.\tag{59}
\]

由于自同态 \(\pi_{\lambda}(g)\) 是酉的，它的逆等于 \(\pi_{\lambda}(g)^*\)，于是

\[
\overline {{\pi_ {i j} ^ {\lambda} (g)}} = \pi_ {j i} ^ {\lambda} (g ^ {- 1}).\tag{60}
\]

于是由 VIII, 第 409 页的公式 (22) 以及第 409 页的 (25) 可知，诸函数 \((d_{\lambda})^{1 / 2}\pi_{ij}^{\lambda}\)（\(\lambda \in \widehat{\mathbf{G}}\)，\(1\leqslant i\leqslant d_{\lambda}\)，\(1\leqslant j\leqslant d_{\lambda}\)）构成希尔伯特空间 \(\mathbf{C}[\mathrm{G}]\) 的一组规范正交基。*

